\section{Conclusions}\label{sec:conclusions}

\ph{Bullet-style conclusions to expand, one per technique, then the comparison.
(1) Dynamic programming gives the exact answer but its table is the bottleneck; divide and conquer did not change the time bound, it changed what fits in memory, and that is what made the full-genome comparison possible.
(2) The linear-time string algorithms only pay off when the pattern is repetitive; on real DNA brute force is already linear in practice, so the theoretical worst case is what justifies KMP, not the average run.
(3) Greedy is provably optimal for the objective it is given, and the ORF experiment shows that the objective is the whole story: maximum count and maximum coverage are different problems.
(4) Backtracking turns an exponential enumeration into a few dozen nodes because the genome is a very strong constraint; its one apparent failure was the frameshift, a fact about the data.
(5) Branch and bound is only as good as its bound: the same optimum, far fewer nodes, but no time saved because the bound is loose where the work is.
(6) Manacher is the only technique whose advantage was invisible on the real data and visible only under an adversarial input.
Comparison across techniques: which scaled best (linear string algorithms and greedy), which needed the most resources (LCS, quadratic, and exhaustive baselines), whether datasets changed the behaviour (yes: real DNA versus adversarial input flips the naive-vs-KMP and centre-vs-Manacher comparisons; the threshold flips the greedy result).}

\ph{Closing paragraph: what was learned about relating experiment to theory: log-log slopes reproduce the exponents, constants matter (numpy versus Python, $2nm$ versus $nm$), and the baseline must be run to know whether the sophisticated algorithm was needed.}
